\PassOptionsToPackage{dvipsnames,svgnames,table}{xcolor}
\documentclass[11pt]{article}
\usepackage[a4paper,margin=22mm]{geometry}
\usepackage[no-math]{fontspec}
\setmainfont{DejaVu Serif}
\usepackage{amsmath,amssymb,amsthm,mathtools,graphicx,xcolor,hyperref,xurl,xspace,ifthen,enumerate}
\usepackage[shortlabels]{enumitem}
\setlength{\emergencystretch}{4em}
\AtBeginDocument{\special{pdf:mapline pzdr Dingbats <uzdr.pfb}}
\SetMathAlphabet{\mathbf}{normal}{OT1}{cmr}{bx}{n}
\SetMathAlphabet{\mathbf}{bold}{OT1}{cmr}{bx}{n}
\newtheorem{theo}{Theorem}[section]
\newtheorem{lemm}[theo]{Lemma}
\newtheorem{thm}[theo]{Theorem}
\newtheorem{proposition}[theo]{Proposition}
\newtheorem{corollary}[theo]{Corollary}
\newtheorem{definition}[theo]{Definition}
\newtheorem{theorem}[theo]{Theorem}
\newtheorem*{remas*}{Remarks}
\newtheorem{remas}[theo]{Remarks}
\newtheorem*{exam*}{Example}
\newtheorem{ques}[theo]{Question}
\providecommand{\og}{«\nobreakspace}
\providecommand{\fg}{\nobreakspace»}
\newtheorem{lem}[theo]{Lemma}
\newtheorem{lemma}[theo]{Lemma}
\newtheorem{prop}[theo]{Proposition}
\newtheorem{coro}[theo]{Corollary}
\newtheorem{corol}[theo]{Corollary}
\newtheorem*{conj*}{Conjecture}
\newtheorem{conj}[theo]{Conjecture}
\newtheorem{conjecture}[theo]{Conjecture}
\newtheorem{Proposition}[theo]{Proposition}
\newtheorem{theoreme}[theo]{Theorem}
\theoremstyle{definition}
\newtheorem{defi}[theo]{Definition}
\newtheorem{exam}[theo]{Example}
\newtheorem{exem}[theo]{Example}
\newtheorem{remark}[theo]{Remark}
\newtheorem{example}[theo]{Example}
\newtheorem{exams}[theo]{Examples}
\newtheorem{rema}[theo]{Remark}
\newtheorem{nota}[theo]{Notation}
\newtheorem{notation}[theo]{Notation}
\newtheorem*{notas}{Notations}
\newtheorem{result}[theo]{Result}
\newtheorem{question}[theo]{Question}
\newtheorem{prob}[theo]{Problem}
\newtheorem*{rema*}{Remark}
\newtheorem*{defi*}{Definition}
\newtheorem*{theo*}{Theorem}
\newtheorem*{lemm*}{Lemma}
\newtheorem*{prop*}{Proposition}
\newcommand{\enoncename}{}
\newtheorem{corpusenonce}[theo]{\enoncename}
\newenvironment{enonce}[2][]{\renewcommand{\enoncename}{#2}\begin{corpusenonce}}{\end{corpusenonce}}
\newenvironment{enonce*}[2][]{\par\medskip\noindent\textbf{#2.}\itshape}{\par\medskip}
\newenvironment{proofc}[1][\proofname]{\begin{proof}[#1]}{\end{proof}}
\newcommand{\Mc}{\mathcal{M}}
\newcommand{\pointir}{.\ ---\ }
\newcommand{\equalenv}[2]{\expandafter\let\csname#1\expandafter\endcsname\csname#2\endcsname\expandafter\let\csname end#1\expandafter\endcsname\csname end#2\endcsname}
\providecommand{\diagbox}[3][]{\begin{tikzpicture}[x=1em,y=1em,baseline=(current bounding box.center)]\useasboundingbox(0,0) rectangle(3,1.4);\draw(0,1.4)--(3,0);\node at(.5,.3){#2};\node at(2.5,1.1){#3};\end{tikzpicture}}
\providecommand{\eg}{e.g.\ }
\providecommand{\ie}{i.e.\ }
\providecommand{\cf}{cf.\ }
\providecommand{\longthanks}[1]{\par\smallskip\noindent #1\par\smallskip}
\providecommand{\qbinom}[2]{\genfrac{[}{]}{0pt}{}{#1}{#2}}

\providecommand{\mathof}[1]{\mathrm{#1}}

\numberwithin{equation}{section}


\def\endappendix{}


\usepackage{latexsym,amsmath,amssymb,verbatim,tikz,tikz-cd}
\usepackage{mathtools}
\usepackage{graphicx,xcolor}

\usepackage[english]{babel}
\usepackage{amssymb}
\usepackage{amsmath}
\usepackage{amsthm}
\usepackage{graphicx}
\usepackage{csquotes}
\usepackage{bm}
\usepackage{enumitem}
\usepackage{subcaption}
\usepackage{caption}
\renewcommand{\thesubfigure}{\Alph{subfigure}}
\usepackage{dirtytalk}
\providecommand{\nopunct}{\spacefactor1007\relax}
\newtheorem{problem}[theo]{Problem}

\newcommand{\sym}{\mathrm{Sym}}
\newcommand{\qsym}{\mathrm{QSym}}
\newcommand{\Des}{\operatorname{Des}}
\newcommand{\Asc}{\operatorname{Asc}}
\newcommand{\Syt}{\operatorname{SYT}}
\newcommand{\Short}{\operatorname{Short}}
\newcommand{\Loop}{\operatorname{Loop}}
\newcommand{\Stable}{\operatorname{Stable}}
\newcommand{\len}{\operatorname{len}}
\newcommand{\height}{\operatorname{height}}
\newcommand{\core}{\operatorname{core}}
\newcommand{\res}{\operatorname{res}}
\newcommand{\row}{\operatorname{row}}
\newcommand{\insrt}{\operatorname{insert}}
\newcommand{\UP}{\operatorname{UP}}
\newcommand{\DOWN}{\operatorname{DOWN}}


\newcommand{\ZZ}{\mathbb{Z}}
\newcommand{\NN}{\mathbb{N}}
\newcommand{\RR}{\mathbb{R}}
\newcommand{\QQ}{\mathbb{Q}}
\newcommand{\PP}{\mathbb{P}}

\newcommand{\F}{\mathcal{F}}
\newcommand{\Q}{\mathcal{Q}}
\newcommand{\A}{\mathcal{A}}
\newcommand{\M}{\mathcal{M}}
\newcommand{\I}{\mathcal{I}}
\newcommand{\PM}{\mathcal{PM}}
\newcommand{\Pcal}{\mathcal{P}}


\newcommand{\odds}{\mathrm{odds}}

\makeatletter
\newcommand{\proofstep}[1]{
  \par
  \addvspace{\medskipamount}
  \textbf{#1\@addpunct{.}}\enspace\ignorespaces
}

\newcommand\restr[2]{{
  \left.\kern-\nulldelimiterspace 
  #1 
  \vphantom{\big|} 
  \right|_{#2} 
  }}
\makeatother












\title{Two sequences of insertions of the same length give Knuth equivalent matchings}
\author{Avichai Marmor}
\date{doi:10.5802/alco.351; CC BY 4.0}
\begin{document}
\maketitle
\noindent\textit{Editorial note.} The counted statement is the author's labelled lemma environment at publisher-source.tex lines 1179-1182; its proof environment opens at line 1184, the line after the blank separator, so the pairing is by adjacency and the retained proof range 1184-1188 is the whole environment.\par\medskip
\noindent\textit{Editorial note.} The two lemmas lem:insert i equivalent insert j (source lines 1157-1160) and lem:insert same chord to equivalent remains equivalent (source lines 1167-1170) are retained as context because the counted proof cites them at source line 1187; they are the only outside labels of the retained proof.\par\medskip
\noindent\textit{Editorial note.} This article carries no label-number dump and the retained counter values are left empty for the whole source, so the printed numbering of the card is the wrapper's default; the identity of the counted item is its author label lem:insert many chords to matching is equivalent.\par\medskip
\noindent\textit{Editorial scope.} One original result of the pinned publisher article, counted once: Knuth equivalence of two iterated insertions of the same length (author label lem:insert many chords to matching is equivalent).\par\medskip
\noindent\textit{Allowed standard prerequisites (editorial).}\begin{itemize}
\item Iterated insertion of short chords. insrt\_\{i\_1\} ... insrt\_\{i\_k\}(m) is the matching obtained by performing the insertions successively at the listed positions.. It is the object compared in the statement.
\item Knuth equivalence. Two matchings are Knuth equivalent when a sequence of elementary Knuth transformations joins them.. It is the conclusion of the statement.
\item The two preceding insertion lemmas. The insertion position may be changed freely, and inserting at the same position preserves Knuth equivalence.. They are the two ingredients of the induction step.
\end{itemize}\medskip
\begin{lemma}
\label{lem:insert i equivalent insert j}
    Let $m \in \M_N$ be a matching and let $1 \le i,j \le N+1$ be indices. Then $\insrt_i(m)$ is Knuth equivalent to $\insrt_j(m)$.
\end{lemma}

\begin{lemma}
\label{lem:insert same chord to equivalent remains equivalent}
    Let $m_1,m_2 \in \M_N$ be Knuth equivalent matchings, and let $1 \le i \le N+1$ be an index. Then $\insrt_i(m_1)$ and $\insrt_i(m_2)$ are Knuth equivalent.
\end{lemma}

\begin{lemma}
\label{lem:insert many chords to matching is equivalent}
    Let $m \in \M_n$ be a matching, and let $i_1,\dots, i_k$ and $j_1,\dots, j_k$ be two sequences. Then the matchings $\insrt_{i_1}\cdots\insrt_{i_k}(m)$ and $\insrt_{j_1}\cdots\insrt_{j_k}(m)$ are Knuth equivalent.
\end{lemma}

\begin{proof}
    We prove the statement by induction on $k$. If $k=0$ then the statement is obvious.

    Next, we assume that $k > 0$, and denote $m_i = \insrt_{i_2}\cdots\insrt_{i_k}(m)$ and $m_j = \insrt_{j_2}\cdots\insrt_{j_k}(m)$. We aim to prove that the matchings $\insrt_{i_1}(m_i)$ and $\insrt_{j_1}(m_j)$ are Knuth equivalent. By the induction hypothesis, $m_i$ and $m_j$ are equivalent. Therefore, by Lemma~\ref{lem:insert same chord to equivalent remains equivalent}, the matchings $\insrt_{i_1}(m_i)$ and $\insrt_{i_1}(m_j)$ are equivalent too. In addition, by Lemma~\ref{lem:insert i equivalent insert j}, the matchings $\insrt_{i_1}(m_j)$ and $\insrt_{j_1}(m_j)$ are equivalent. Therefore, the matchings $\insrt_{i_1}(m_i)$ and $\insrt_{j_1}(m_j)$ are equivalent.
\end{proof}

\end{document}
